\documentclass[10pt,a4paper]{amsart}
\usepackage[margin=19mm]{geometry}
\usepackage{amsmath,amssymb,mathtools}
\usepackage[scr=rsfs,cal=euler]{mathalfa}
\usepackage{xfrac}
\usepackage[unicode,hidelinks,bookmarks=false]{hyperref}
\newcommand{\ie}{i.e.\ }
\newcommand{\cf}{cf.\ }
\newcommand{\resp}{respectively\ }
\newcommand{\Subsection}{\subsection}
\providecommand{\nobreakdash}{\nobreak\hbox{-}}
\setlength{\emergencystretch}{2em}
\multlinegap0pt
\usepackage{amssymb}
\newtheorem{lemma}{Lemma}
\newtheorem{theorem}{Theorem}
\DeclareMathOperator{\dom}{dom}
\DeclareMathOperator{\girth}{girth}
\newcommand{\sep}{\;|\;}
\newcommand{\seq}{\subseteq}
\title{From $b$-Coloring to $b^*$-Coloring: Large Girth and Parameterized Complexity}
\author{Jakub Balab\'an, Oliver Bukor}
\date{Source: arXiv:2608.03819v2 (CC BY 4.0)}
\begin{document}
\maketitle

\noindent\emph{Source and licence.} From $b$-Coloring to $b^*$-Coloring: Large Girth and Parameterized Complexity. Version arXiv:2608.03819v2 (CC BY 4.0), distributed by arXiv under the Creative Commons Attribution 4.0 International licence, http://creativecommons.org/licenses/by/4.0/, as recorded in the arXiv licence metadata for this exact version and shown on the abstract page https://arxiv.org/abs/2608.03819v2. Full licence notice: \texttt{licenses/arxiv-2608.03819-CC-BY-4.0.txt}.\par\smallskip

\noindent\textit{Editorial scope.} Counted unit: the article's theorem thm:girth7 (source0.tex lines 438--440). The article's complete original proof of it is delivered with it (source0.tex lines 441--474, one \texttt{proof} environment of 34 lines). No mathematics is written, repaired or re-derived: the statement and the proof are the article's own lines, in the article's own notation, and no retained line is substituted or re-flowed.  Retained uncounted as the source-local context the proof needs: lines 405--407; lines 423--427.  Nothing was omitted from any retained span, so the package carries no omission record.  No retained line contains a citation, so the wrapper carries no bibliography and no bibliography entry.  No retained line contains a figure environment, an image-inclusion command or drawing code, so no image is delivered and the package declares no asset dependency.  Editorial formatting only, in addition: an AMS \texttt{amsart} preamble that restates the article's own declarations and macros used by the retained lines, namely the environments \texttt{lemma}, \texttt{theorem} on separate counters, together with the standard packages those lines need.  The retained environments print in this document as Lemma 1, Theorem 1 in order of appearance, whereas the article prints the same items with the numbering of its own full document.  The complete native source of the article as distributed in the arXiv e-print for this exact version is bundled unchanged as \texttt{sources/206787/source0.tex}.
\begin{lemma}\label{lem:g6-init}
If $G$ is a graph of girth at least 6, then there is a partial $b^*$-coloring $\chi$ of $G$ using $k := m^*(G)+1$ colors such that $|\dom(\chi)| = k^2 -2k + 2$.
\end{lemma}

\begin{lemma}\label{lem:recolor}
Let $G$ be a graph, let $\chi$ be a partial $k$-coloring of $G$, let $(u, U)$ be a $\chi$-$b^*$-core, and let $v \in \dom(\chi) \setminus U$.
If $v$ is not a $b$-vertex in $\chi$, 
then there is a color $c \in [k] \setminus \{\chi(v)\}$ such that $u$ is a $b^*$-vertex in $\chi' := \chi[v \mapsto c]$ and $\chi'$ is proper.
\end{lemma}

\begin{theorem}\label{thm:girth7}
If $G$ is a graph of girth at least $7$, then $b^*(G) = m^*(G) + 1$ and $G$ is $b^*$-monotonic.
\end{theorem}

\begin{proof}
Let $G$ be an $n$-vertex graph, let $k = m^*(G) + 1$, and let $i_0 = k^2 - 2k + 2$.
For each $i \in [i_0, n]$, we will construct a partial coloring satisfying the following invariant: $\chi_i$ is a $k$-$b^*$-coloring and $|\dom(\chi_i)| = i$. 
Then, the total coloring $\chi_n$ will certify that $b^*(G) = m^*(G) + 1$.
Let $\chi_{i_0}$ be the coloring provided by Lemma~\ref{lem:g6-init} and observe that it satisfies the invariant.

Suppose that $\chi_i$ is defined for some $i \in [i_0, n-1]$ and let $v \in \overline{\dom(\chi_i)}$.
For each $c \in [k]$, let $V_c = N(v) \cap \chi_i^{-1}(c)$.
First, suppose that for each $c \in [k]$, there is a $\chi_i$-$b$-vertex $w \in V_c$.
Observe that $w$ has degree at least $k$; it has $k-1$ neighbors colored by $\chi_i$ plus $v$.
Hence, $m^*(v) \ge k$, which is a contradiction with $m^*(G) = k-1$.
Hence, there is a color $c \in [k]$ such that $V_c$ contains no $\chi_i$-$b$-vertices; let us fix such a color $c$. Note that $V_c = \emptyset$ is possible.

Let $(u, U)$ be a $\chi_i$-$b^*$-core.
If $V_c \cap U = \emptyset$, then we successively apply Lemma~\ref{lem:recolor} to vertices in $V_c$ to obtain a proper coloring $\chi'$ such that $u$ is a $b^*$-vertex in $\chi'$ and $c \notin \chi'(N(v))$.
Note that after one application of the lemma, no other vertex in $V_c$ becomes a $b$-vertex because $V_c$ is an independent set in $G$, so the construction of $\chi'$ is valid.
Now the coloring $\chi_{i+1} := \chi' \cup \{(v, c)\}$ satisfies the invariant.
Hence, we may assume that $V_c \cap U \ne \emptyset$.

Suppose that there is a color $d \in [k] \setminus \{c\}$ such that $V_d$ contains no $\chi_i$-$b$-vertices.
Since $G[U]$ has diameter at most 4, $G$ has girth at least 7 and $V_c \cap U \ne \emptyset$, we have $V_d \cap U = \emptyset$.
Hence, we may, analogously to the previous paragraph, use Lemma~\ref{lem:recolor} to obtain a proper coloring $\chi'$ such that $u$ is a $b^*$-vertex in $\chi'$ and $d \notin \chi'(N(v))$, and the coloring $\chi_{i+1} = \chi' \cup \{(v, d)\}$ satisfies the invariant.

Finally, suppose that no such color $d$ exists, i.e., for each color $d \in [k] \setminus \{c\}$, there is a $\chi_i$-$b$-vertex $v_d \in V_d$.
Recall that by choice of $c$, $V_c$ contains no $\chi_i$-$b$-vertices.
Hence, for each $w \in V_c$, there is a color $c_w \in [k]$ such that $\chi_i[w \mapsto c_w]$ is proper.
Since $V_c$ is an independent set, the coloring $\chi' := \chi_i \setminus \{(w, c) \sep w \in V_c\} \cup \{(w, c_w) \sep w \in V_c\}$ is also proper.
Let us define $\chi_{i+1} = \chi' \cup \{(v, c)\}$, and observe that $\chi_{i+1}$ is proper.
Since $N(v)$ is an independent set in $G$, $v_d$ is a $b$-vertex colored with $d$ also in $\chi_{i+1}$, for every color $d \in [k] \setminus \{c\}$.
Hence, $v$ is a $b^*$-vertex in $\chi_{i+1}$, which means that $\chi_{i+1}$ satisfies the invariant.

We have proven that $b^*(G) = m^*(G) + 1$.
To show that $G$ is $b^*$-monotonic, it suffices to observe that $m^*(H_1) \le m^*(H_2)$ and $\girth(H_1) \le \girth(H_2)$ whenever $H_1$ and $H_2$ are two graphs satisfying $H_1 \seq_i H_2$.
\qed\end{proof}

\end{document}
